\section{Encuestas}

\subsection{Diseño de encuesta}
\subsubsection{Objetivo de medición}
\noindent El objetivo de la encuesta es medir la satisfacción del turista al finalizar el servicio, identificar oportunidades de mejora y fortalecer la calidad de atención de Machu Picchu Exclusive Tours.
\subsubsection{Estructura de preguntas}
\noindent Se recomienda una estructura breve y clara, combinando:
\begin{itemize}
    \item preguntas de valoración (por ejemplo, atención, puntualidad y experiencia general),
    \item una pregunta abierta para comentarios del cliente, y
    \item una pregunta final de recomendación del servicio.
\end{itemize}

\subsection{Aplicación}
\subsubsection{Momento de envío}
\noindent El envío debe realizarse cuando el cliente culmina su tour o dentro de las primeras 24 horas posteriores, para obtener respuestas más precisas y útiles.
\subsubsection{Canal de recolección}
\noindent El canal principal de recolección será el correo electrónico, según el medio de contacto registrado del pasajero. Lo importante es mantener consistencia en el canal por tipo de cliente.

\subsection{Análisis}
\subsubsection{Lectura de resultados}
\noindent Revise periódicamente los resultados para detectar patrones: puntos fuertes del servicio, incidencias repetitivas y sugerencias frecuentes de mejora.
\subsubsection{Acciones de mejora}
\noindent Con base en los resultados, defina acciones concretas de mejora operativa y comercial. Se recomienda documentar cada ajuste para evaluar su impacto en las siguientes mediciones.

\subsection{Apoyo visual recomendado}
\noindent Para esta sección, la parte visual del flujo se aprende mejor con una demostración en video. Como material de apoyo oficial, utilice el siguiente enlace:

\noindent \href{https://youtu.be/Iv9yd-VFe90?si=qw1tjA2YIO08mhW3}{Video de apoyo: configuración y uso de encuestas}
